\documentclass[11pt]{exam}
\usepackage[a4paper,margin={.1\paperheight,.1\paperwidth},marginratio=1:1]{geometry}
\usepackage[french]{babel}
\usepackage[T1]{fontenc}

% Choix des options et des infos d’en-tête
\usepackage[ds,prof]{profnsi}
\profnsisetup{
	niveau=Première,
	matiere=NSI,
	titre=Représentation des données,
	chapitre=3, 
}


\begin{document}
	\creerDoc
	
	\partie{Conversions}
	
	\exercice{Convertissez les nombres binaires suivants :}
	
	\begin{enumerate}
		\item $101011_2$ = \correction{$43_{10}$}
		\item \( 101011_2 \) = 
		\item \( 101011_2 \) = 
		\item \( 101011_2 \) = 
		\item \( 101011_2 \) = 
	\end{enumerate}
	


	\exercice{Décimal vers Binaire}
	
	\begin{enumerate}
		\item \( 101011_2 \) = 
		\item \( 101011_2 \) = 
		\item \( 101011_2 \) = 
		\item \( 101011_2 \) = 
		\item \( 101011_2 \) = 
	\end{enumerate}

	
	\exercice{Binaire vers hexadécimal}
	
		\partie{Conversions}
	
	\exercice{Convertissez les nombres binaires suivants :}
	
	\begin{enumerate}
		\item \( 101011_2 \) = 
		\item \( 101011_2 \) = 
		\item \( 101011_2 \) = 
		\item \( 101011_2 \) = 
		\item \( 101011_2 \) = 
	\end{enumerate}
	
	
	
	\exercice{Décimal vers Binaire}
	
	\begin{enumerate}
		\item \( 101011_2 \) = 
		\item \( 101011_2 \) = 
		\item \( 101011_2 \) = 
		\item \( 101011_2 \) = 
		\item \( 101011_2 \) = 
	\end{enumerate}
	
	
	\exercice{Binaire vers hexadécimal}
	
	
\end{document}
